\subsubsection{仅加边动态 scc}
每条边的生效时间是不小于他的加入时间的，整体二分这个时刻，维护 $(l,r,E)$ 表示生效时间在 $[l,r]$ 内的边。加入 $[l,mid]$ 中的边考察 scc 变化情况判断生效时间在左区间还是右区间（也可能不生效）。

\begin{lstlisting}
int dfn[100010],low[100010],tim;
int st[100010],tp;
bool in[100010];
int bel[100010],scc;
vector<int>e[100010];
vector<int>node;
void dfs(int u){
    dfn[u]=low[u]=++tim;
    st[++tp]=u,in[u]=1;
    for(auto v:e[u]){
        if(!dfn[v])dfs(v),low[u]=min(low[u],low[v]);
        else if(in[v])low[u]=min(low[u],dfn[v]);
    }if(dfn[u]==low[u]){
        scc++;int now;while(1){
            now=st[tp--];
            bel[now]=scc;in[now]=0;
            if(now==u)break;
        }
    }
}
void clear(){
    tim=0,tp=0;scc=0;
    for(auto u:node)dfn[u]=low[u]=0,bel[u]=0,e[u].clear();
    node.clear();
}
int U[100010],V[100010];
vector<int>con[100010];
int fa[100010];
int find(int x){
    if(fa[x]==x)return x;
    return fa[x]=find(fa[x]);
}
void un(int u,int v){
    if(find(u)==find(v))return;
    u=find(u),v=find(v);
    fa[u]=v;
}
void work(int l,int r,vector<array<int,3>>&E){// E 中储存 (u,v,i)，表示在加入全部 $[1,l-1]$ 的边后，u 和 v 所在的 scc 编号。这样接着在这个图上做可以简化。
    // printf("Working on [%d, %d]\n",l,r);
    // for(auto [u,v,i]:E)printf("(%d, %d, %d) ",u,v,i);cout<<endl;
    int mid=(l+r)>>1;
    vector<array<int,3>> e1,e2;e1.clear(),e2.clear();
    for(auto [u,v,i]:E){
        node.pb(u),node.pb(v);
        if(i<=mid)e[u].pb(v);
    }
    for(auto u:node)if(!dfn[u])dfs(u);
    if(l==r){
        for(auto [u,v,i]:E){
            if(bel[u]==bel[v])un(U[i],V[i]);//有了“生效”状态后即可转化为并查集维护。
        }clear();
        return;
    }
    for(auto [u,v,i]:E){
        if(bel[u]!=bel[v])e2.pb({bel[u],bel[v],i});
        else if(i<=mid&&bel[u]==bel[v])e1.pb({u,v,i});
    }clear();
    work(l,mid,e1),work(mid+1,r,e2);
}
\end{lstlisting}

